\documentclass[10pt,twocolumn]{article}
\usepackage[scaled]{helvet}
\renewcommand\familydefault{\sfdefault}
\usepackage[T1]{fontenc}
\usepackage[margin=0.7in]{geometry}
\usepackage{titlesec}
\usepackage{enumitem}
\usepackage{parskip}
\setlength{\columnsep}{0.24in}
\setlist[itemize]{leftmargin=1.2em, topsep=0.2em, itemsep=0.18em}
\titleformat{\section}{\large\bfseries}{\thesection}{1em}{}

\begin{document}
\title{\vspace{-1.6cm}AWS SageMaker Deployment}
\author{AWS ML Study Notes}
\date{}
\maketitle
\small

\section*{Overview}
SageMaker Deployment turns trained models into managed inference endpoints or batch inference jobs. It packages model artifacts, container logic, scaling behavior, and network settings into a serving configuration.

The value is operational consistency. Teams can promote models through controlled environments instead of treating deployment as a manual afterthought.

\section*{Core Concepts}
\begin{itemize}
\item A deployment usually starts with a model definition that points to the serving image and the model artifact in S3.
\item Endpoint configurations describe instance types, variant weights, autoscaling behavior, or specialized modes such as asynchronous or serverless inference.
\item SageMaker supports patterns such as multi model endpoints, blue green style updates through variants, and data capture for later monitoring.
\end{itemize}

\section*{How It Works}
\begin{itemize}
\item After a model is approved, the deployment system creates or updates the endpoint configuration and then rolls traffic to the new endpoint variant.
\item Clients send inference requests to the hosted endpoint, while logs, metrics, and optionally request payload capture flow to surrounding AWS services.
\item For offline scoring, the same model package can be used in batch transform rather than a permanently running real time endpoint.
\end{itemize}

\section*{AI and ML Relevance}
\begin{itemize}
\item Deployment design affects latency, throughput, cost, and release safety. The right mode depends on traffic profile, payload size, and how quickly results are needed.
\item ML teams also care about safe promotion, shadowing, rollback, and post deployment quality checks because correctness drift can be more subtle than simple uptime failures.
\end{itemize}

\section*{Operational Notes}
\begin{itemize}
\item Use autoscaling carefully. If cold starts or model load time are high, a low minimum instance count can cause unstable latency.
\item Capture inference data where governance allows it so you can debug failure cases, audit outcomes, and build monitoring or retraining signals.
\item Separate release approval from endpoint mutation. Production deployment should be traceable to a specific model version and image version.
\end{itemize}

\section*{What To Remember}
\begin{itemize}
\item Training creates a model artifact; deployment turns it into a callable service contract.
\item Real time, asynchronous, serverless, and batch inference exist for different traffic and cost profiles.
\item The hardest part is rarely the API call. It is the release discipline around versions, rollbacks, and observability.
\end{itemize}

\end{document}
